% !TeX root = ../tango.tex
\section{Implementation of TANGO}
\label{sec:implementation-of-tango}

\tango{} is implemented on top of vLLM 0.17.0 by extending its request admission path and scheduler.
The dispatcher and allocator run in the CPU-side scheduling loop and produce per-request token allocations for each engine step.
The resulting batch is passed to vLLM through its existing scheduler--executor interface, reusing continuous batching, chunked prefill, PagedAttention-based KV-cache management, tensor parallelism, sampling, and streaming output.

\textbf{SLO Metadata and Runtime State:}
We extend request admission to carry two optional SLO fields, TTFT and TPOT, in milliseconds.
A missing SLO is treated as an absent constraint, allowing best-effort requests to use the same serving path.
Alongside these fields, the scheduler maintains each request's arrival timestamp, prompt length, computed-token progress, output-token count, and latest output timestamp in host memory.
These fields are updated together with vLLM's request state after each engine step.
The first emitted token activates the TPOT objective, and subsequent token timestamps advance the next-token deadline.

\textbf{Queue and Preemption Hooks:}
We add queue and preemption hooks to vLLM's waiting-request admission path.
The queue hook applies the phase-aware priority and tie-breaking rules in Section~\ref{sec:deadline-aware-request-dispatcher}.
On KV-cache allocation failure, the preemption hook selects the running request with the latest deadline and requires the waiting request's deadline to be strictly earlier.
If permitted, vLLM reclaims the victim's KV blocks, moves it from \texttt{RUNNING} to \texttt{WAITING}, and retries allocation.
The victim keeps its SLOs, output progress, and the timestamps used to construct its deadline.

\textbf{Per-Step Token Allocation:}
We fit the step-time model coefficients from offline profiling for the deployed model and hardware configuration.
At each iteration boundary, the allocator summarizes the active decode workload and collects its next-token deadlines from scheduler-maintained state.
It then establishes the time budget using the percentile policy in Section~\ref{sec:step-time-model-guided-token-allocator}.
As the scheduler processes ongoing prefill requests and admits waiting requests, their chunks share this budget and the engine's maximum per-step token capacity.
The allocator updates the cumulative prefill token count and attention work after each allocation, using the candidate's chunk length and already-computed context length.
A chunk that exceeds either limit is truncated to its largest feasible size; unallocated tokens remain for later steps.
The resulting token allocations are forwarded to the executor through the standard scheduling output.
